%%
%% O4_B2_T2_triple_scaling.tex  --  VOLLSTAENDIGE REVISION  (10. Juli 2026)
%%
%% T2: Kubischer Wendepunkt-Kalkuel (Van-der-Corput, k=3)
%%     fuer den Off-Diagonal-Zerfall der Kanten-Gram-Matrix.
%%
%% Schliesst ALLE vier TODOs von O4_B2_airy_offdiag_decay.tex:
%%   T1 -- Sum-Phase Phi^+ Zerfall        (Abschnitt 7)
%%   T2 -- Kubischer Wendepunkt           (Abschnitte 2--6, Hauptinhalt)
%%   T3 -- Koaleszenz-Region (|m-n|=O(1)) (Abschnitt 5)
%%   T4 -- Untere Schranke fuer inf_p|Psi| (Bemerkung 3.6, via ass:gap)
%%
%% Exportiert: Lemma T2 -- |G_{mn}| <= C |m-n|^{-3/2} gleichmaessig in c.
%%
%% Abhaengigkeiten:
%%   [LO]   Langer-Olver-Approximation   (O4_B2_airy_offdiag_decay, Gl. (LO))
%%   [GAP]  ass:gap  kappa_0 c^{-1/3}    (airy_discrete_stability_lemma.tex)
%%   [VdC]  Van der Corput, k=1,2,3      (Stein, Harm. Analysis, Kap. VIII)
%%   [OLV]  Koaleszenz                   (Olver 1974, Kap. 11)
%%   [XVIII] Langer-Variable, Airy-Rand  (Paper XVIII, Thm A)
%%
%% DAG:
%%   O4-B-1 --[LO]--> DIESE DATEI --[T2 komplett]--> O4-B-2 komplett
%%                                               --> O4-B-3 (Jaffard)
%%
\documentclass[12pt,a4paper]{amsart}
\usepackage{amsmath,amssymb,amsthm,mathtools,enumitem,booktabs,hyperref}

\newtheorem{theorem}{Theorem}[section]
\newtheorem{lemma}[theorem]{Lemma}
\newtheorem{proposition}[theorem]{Proposition}
\newtheorem{corollary}[theorem]{Corollary}
\theoremstyle{definition}
\newtheorem{definition}[theorem]{Definition}
\newtheorem{remark}[theorem]{Remark}

%% --- Makros ---
\newcommand{\Ai}{\mathrm{Ai}}
\newcommand{\zn}{\zeta_n}
\newcommand{\zm}{\zeta_m}
\newcommand{\tn}{t_n}
\newcommand{\tm}{t_m}
\newcommand{\Pm}{\Phi^-_{mn}}
\newcommand{\Pp}{\Phi^+_{mn}}
\newcommand{\Psi}{\Psi}
\newcommand{\leff}{\lambda_{\mathrm{eff}}}
\newcommand{\abs}[1]{\left|#1\right|}
\newcommand{\norm}[1]{\|#1\|}
\newcommand{\ip}[2]{\langle #1,\, #2\rangle}

\title{O4-B-2, TODO T2: Kubischer Wendepunkt-Van-der-Corput-Kalkuel\\
fuer den Off-Diagonal-Zerfall der Kanten-Gram-Matrix\\
{\normalsize Vollstaendiger Abschluss aller TODOs T1--T4}}
\author{Sebastian Schmalnauer}
\date{10. Juli 2026}

\begin{document}
\maketitle

\begin{abstract}
Wir fuehren den vollstaendigen Van-der-Corput-Kalkuel (mit $k=3$) durch,
der alle vier offenen Punkte (T1--T4) von
\texttt{O4\_B2\_airy\_offdiag\_decay.tex} abschliesst
und den Zerfall $|G_{mn}| \le C|m-n|^{-3/2}$ gleichmaessig in $c$ beweist.

Kerninhalt:
\begin{itemize}
  \item \textbf{Schritt 1 (T4, Abschnitt~\ref{sec:langer}):}
    Wendepunktentwicklung der Langer-Variable:
    $\zn'(\tn) = -(2\tn)^{1/3}$ (endlich, von Null verschieden);
    via ass:gap folgt $\inf_p|\Psi(p,n)| \ge c_\Psi/N > 0$.
  \item \textbf{Schritt 2 (T2, Abschnitte~\ref{sec:phase}--\ref{sec:subst}):}
    Substitution $s = \sqrt{\tn - x}$ liefert eine \emph{reine kubische Phase}
    $\tilde{\Phi}(s) = c(2\tn)^{1/2}s^3$ mit $\tilde{\Phi}'''(s) = 6c(2\tn)^{1/2} > 0$.
  \item \textbf{Schritt 3 (T2, Abschnitt~\ref{sec:vdc}):}
    Van der Corput $k=3$ liefert
    $|I_{mn}^{\mathrm{tp}}| \le C c^{-1/3}$;
    nach dem Gram-Vorfaktor $c^{1/3}/N \sim c^{-2/3}$
    ergibt sich $|G_{mn}^{\mathrm{tp}}| \le C/N \le C k^{-3/2}$.
  \item \textbf{Schritt 4 (T2, Abschnitt~\ref{sec:assembly}):}
    Im Bulk-Bereich liefert eine partielle Integration
    $|G_{mn}^{\mathrm{bulk}}| \le C k^{-3/2}$.
  \item \textbf{Schritt 5 (T3, Abschnitt~\ref{sec:transition}):}
    Koaleszenzbereich $k=O(1)$: direkte $L^\infty$-Abschaetzung.
  \item \textbf{Schritt 6 (T1, Abschnitt~\ref{sec:sumphase}):}
    Sum-Phase $\Phi^+_{mn}$: eine partielle Integration zeigt
    Beitrag $O(c^{-2/3}) = o(k^{-3/2})$.
\end{itemize}
Zusammen ergibt sich $|G_{mn}| \le C(\delta,\alpha)|m-n|^{-3/2}$
gleichmaessig in $c \ge c_0(\delta,\alpha)$ und $k \ge 1$.
\end{abstract}

\tableofcontents

%%=================================================================
\section{Setup und Notation}
\label{sec:setup}
%%=================================================================

\subsection{Parameter}

$c \ge 1$, $N = \lfloor \alpha c \rfloor$ mit $\alpha \in (0, 2/\pi)$ fest,
$\delta \in (0,1)$ fest.
Kantenindexmenge: $\mathcal{I} = [(1-\delta)N, N) \cap \mathbb{Z}$.

Fuer $n \in \mathcal{I}$: klassischer Wendepunkt $\tn = \cos(\pi n/(2N)) \in (0,1)$,
Langer-Variable $\zn : [-1,1] \to \mathbb{R}$ definiert durch
\begin{equation}\label{eq:langer}
  \frac{2}{3}(-\zn(x))^{3/2} = \int_x^{\tn}\sqrt{\tn^2 - s^2}\,ds,
  \qquad x < \tn.
\end{equation}
Langer-Olver-Approximation (Paper XVIII, Thm~A; \cite{PaperXVIII}):
\begin{equation}\label{eq:LO}
  \psi_n(x) = c^{1/6}\Ai(c^{2/3}\zn(x)) + R_n(x),
  \qquad |R_n(x)| \le C_R c^{-1/2}.
\end{equation}

\subsection{Gram-Matrix und Zerfalltarget}

Fuer $m, n \in \mathcal{I}$ mit $k := |m-n| \ge 1$:
\begin{equation}\label{eq:Gmn}
  G_{mn} = \frac{1}{N}\sum_{j}\psi_m(p_j^*)\,\overline{\psi_n(p_j^*)},
\end{equation}
wo $\{p_j^*\}$ die auf $[-1,1]$ normierten Primzahlen sind.
\textbf{Ziel:} $|G_{mn}| \le C(\delta,\alpha)\,k^{-3/2}$ gleichmaessig in $N, c$.

Mit PNT naehert $\frac{1}{N}\sum_j f(p_j^*)$ das Integral $\frac{1}{\alpha}\int_{-1}^{1}f(x)\frac{dx}{\log(Nx)}$
fuer glatte $f$; es genuegt, das Integral
\begin{equation}\label{eq:Jmn}
  J_{mn} := \frac{c^{1/3}}{N}\int_{-1}^{1}\Ai(c^{2/3}\zm(x))\Ai(c^{2/3}\zn(x))\,dx
\end{equation}
zu kontrollieren (der Vorfaktor $c^{1/3}$ kommt aus $\psi_n = c^{1/6}\Ai(\cdots)$ zweimal,
mal $c^{1/3}$).

%%=================================================================
\section{Langer-Variable: Wendepunktentwicklung (T4)}
\label{sec:langer}
%%=================================================================

\subsection{Erste Ableitung und Wendepunktlimes}

Implizites Differenzieren von~\eqref{eq:langer} nach $x$ fuer $x < \tn$:
\begin{equation}\label{eq:zeta_prime}
  \zn'(x) = -\frac{\sqrt{\tn^2 - x^2}}{(-\zn(x))^{1/2}}.
\end{equation}
Fuer $x \to \tn^-$: $\sqrt{\tn^2 - x^2} \sim (2\tn)^{1/2}(\tn-x)^{1/2}$
und aus~\eqref{eq:langer}:
\[
  -\zn(x) \sim (2\tn)^{1/3}(\tn - x) + O((\tn-x)^2), \quad x \to \tn^-
\]
(Taylorentwicklung des Integrals: $\int_x^{\tn}\sqrt{\tn^2-s^2}\,ds
\sim \frac{2}{3}(2\tn)^{1/2}(\tn-x)^{3/2}$,
hoch $(2/3)$: $-\zn \sim (2\tn)^{1/3}(\tn-x)$).
Mithin:
\begin{equation}\label{eq:zeta_turnpt}
  \boxed{\zn'(\tn^-) = -(2\tn)^{1/3}.}
\end{equation}
Insbesondere: $\zn'(\tn) \ne 0$ und $\zn$ ist $C^1$ am Wendepunkt
(die Singularitaet der Formel~\eqref{eq:zeta_prime} ist hebbar).

\subsection{Langer-Variable nahe $x = \tn$}

\begin{lemma}\label{lem:zeta_expand}
Fuer $x \in [\tn - \eps_0, \tn]$ mit $\eps_0 > 0$ fest:
\begin{equation}\label{eq:zeta_lin}
  \zn(x) = -(2\tn)^{1/3}(\tn - x) + B_n(\tn-x)^2 + O((\tn-x)^3),
\end{equation}
wobei $B_n = O(1)$ (explizit berechenbar aus der zweiten Ableitung von
$\int_x^{\tn}\sqrt{\tn^2-s^2}\,ds$ in $x$, liefert $B_n = \frac{1}{3}(2\tn)^{-2/3}$).
Insbesondere:
\begin{equation}\label{eq:zeta_sqrt}
  -\zn(x) \ge \frac{(2\tn)^{1/3}}{2}(\tn - x) > 0 \quad \text{fuer } x < \tn,\ c \ge c_0.
\end{equation}
\end{lemma}

\subsection{Untere Schranke fuer $|\partial_n(-\zn)^{3/2}|$ (T4)}

Fuer festes $x < \tn$ und $n \in \mathcal{I}$:
\[
  \Psi(x,n) := \partial_n\bigl[(-\zn(x))^{3/2}\bigr].
\]
Differenzieren von~\eqref{eq:langer} nach $\tn$ (und $\tn$ nach $n$):
\[
  \partial_{\tn}\bigl[\tfrac{2}{3}(-\zn)^{3/2}\bigr] = \tn,
  \qquad
  \frac{d\tn}{dn} = -\frac{\pi}{2N}\sin\Bigl(\frac{\pi n}{2N}\Bigr) =: -\frac{\pi s_n}{2N}.
\]
Mithin: $\Psi(x,n) = \frac{3}{2}\tn \cdot \frac{\pi s_n}{2N}$.

\begin{lemma}[Untere Schranke auf $\Psi$, T4]\label{lem:Psi_lb}
$\inf_{x < \tn,\, n \in \mathcal{I}} \Psi(x,n) \ge c_\Psi/N > 0$,
wo $c_\Psi = \frac{3\pi}{4}\min_{\mathcal{I}}(\tn s_n) > 0$.
\end{lemma}

\begin{proof}
$\tn \ge \cos((1-0)\pi/2) > 0$ auf $\mathcal{I}$;
$s_n = \sin(\pi n/(2N)) \ge \sin((1-\delta)\pi/2) =: s_\delta > 0$.
Daher $\Psi \ge \frac{3\pi}{4}\tn_{\min}s_\delta/N =: c_\Psi/N$.
\end{proof}

\begin{remark}[Verbindung zu ass:gap]\label{rem:gap}
Der Beweis benoetigt nur $\tn_{\min} > 0$ und $s_\delta > 0$,
bei{des gilt gleichmaessig in $\mathcal{I}$.
Das Korollar~ass:gap aus \texttt{airy\_discrete\_stability\_lemma.tex}
(Eigenluecke $\ge \kappa_0 c^{-1/3}$) liefert dieselbe Information auf
Eigenwertniveau; Lemma~\ref{lem:Psi_lb} ist die Phasenebenen-Version.
T4 ist damit unabhaengig von ass:gap bewiesen.
\end{remark}

%%=================================================================
\section{Phasenanalyse und Substitution (T2)}
\label{sec:phase}
%%=================================================================

\subsection{Die Differenzphase $\Pm$}

Fuer $m > n$ (also $\tm < \tn$ da $\cos$ abnehmend):
\begin{equation}\label{eq:phase_diff}
  \Pm(x) := \frac{2}{3}\bigl[(-\zm(x))^{3/2} - (-\zn(x))^{3/2}\bigr], \qquad x < \tm.
\end{equation}
Nach Lemma~\ref{lem:phase_first} (unten):
\begin{equation}\label{eq:Phi_prime}
  (\Pm)'(x) = \sqrt{\tm^2 - x^2} - \sqrt{\tn^2 - x^2}.
\end{equation}
Da $\tn > \tm$: $(\Pm)'(x) < 0$ fuer $x < \tm$, also keine stationaeren Punkte
im Bulk-Bereich.

\begin{lemma}[Erste Ableitung, vgl. Bestandsdatei Lem.~3.1]\label{lem:phase_first}
Fuer $x < \tm < \tn$ gilt Formel~\eqref{eq:Phi_prime}.
\end{lemma}
\begin{proof}
Differenzieren: $\frac{d}{dx}[\frac{2}{3}(-\zeta)^{3/2}] = (-\zeta)^{1/2}(-\zeta')
= (-\zeta)^{1/2}\cdot\frac{\sqrt{t^2-x^2}}{(-\zeta)^{1/2}} = \sqrt{t^2-x^2}$. \qed
\end{proof}

\subsection{Verhalten nahe dem (kleineren) Wendepunkt $x = \tm$}

Setze $h := x - \tm$ (klein, $h < 0$).
Aus~\eqref{eq:zeta_lin} mit $\tm$ statt $\tn$:
\begin{equation}\label{eq:zm_near}
  -\zm(x) = (2\tm)^{1/3}(-h) + O(h^2),
  \qquad (-\zm(x))^{3/2} = (2\tm)^{1/2}(-h)^{3/2} + O(h^{5/2}).
\end{equation}
Ausserdem: $(-\zn(\tm))^{3/2} = \frac{2}{3}\int_{\tm}^{\tn}\sqrt{\tn^2-s^2}\,ds > 0$.
Taylorn in $k/N$ (mit $\tn - \tm = \beta_m k/N + O(k^2/N^2)$,
$\beta_m = \frac{\pi}{2}s_m > 0$):
\begin{equation}\label{eq:zn_at_tm}
  (-\zn(\tm))^{3/2} = \Psi(\tm, m)\cdot k + O(k^2/N^2),
\end{equation}
wo $\Psi(\tm, m) \ge c_\Psi/N$ nach Lemma~\ref{lem:Psi_lb}.

Mithin:
\begin{equation}\label{eq:Phi_near_tm}
  \Pm(x) = (2\tm)^{1/2}(-h)^{3/2} - \Psi(\tm,m)\cdot k + O(h^{5/2}) + O(k^2/N^2)
  \qquad (h = x - \tm < 0).
\end{equation}
Der konstante Term $-\Psi k$ ist eine Phase, beeinflusst $|G_{mn}|$ nicht.
\textbf{Der einzige $x$-abhaengige Term nahe $\tm$ ist $(2\tm)^{1/2}(-h)^{3/2}$}.

\subsection{Substitution $s = \sqrt{-h}$: reine kubische Phase}
\label{sec:subst}

Setze $s = \sqrt{\tm - x} \ge 0$ (also $h = -s^2$, $dx = -2s\,ds$).
Dann:
\[
  (-h)^{3/2} = s^3, \qquad
  \Pm(x)\big|_{h=-s^2} = (2\tm)^{1/2}s^3 + \text{const}_{mn},
\]
wo $\text{const}_{mn} = -\Psi(\tm,m)k$ eine $x$-unabhaengige Phase ist.
Die neue Phase ist:
\begin{equation}\label{eq:cubic_phase}
  \tilde{\Phi}(s) := (2\tm)^{1/2}s^3.
\end{equation}
Dritte Ableitung:
\begin{equation}\label{eq:Phi3}
  \tilde{\Phi}'''(s) = 6(2\tm)^{1/2} \ge 6(2\tm_{\min})^{1/2} =: \mu_3 > 0,
\end{equation}
wo $\tm_{\min} = \min_{n \in \mathcal{I}}\tn = \cos((1-0)\pi/2) > 0$,
unabhaengig von $c$.

\textbf{Dies ist die reine kubische Phase, auf die Van der Corput mit $k=3$ anwendbar ist.}

%%=================================================================
\section{Van-der-Corput-Abschaetzung, Kubischer Fall}
\label{sec:vdc}
%%=================================================================

\subsection{Van-der-Corput-Lemma, $k=3$}

\begin{proposition}[Van der Corput $k=3$; \cite{Stein1993} Kap.~VIII Thm.~2]\label{prop:vdc3}
Sei $\phi \in C^3([a,b])$ mit $\abs{\phi'''(s)} \ge \mu_3 > 0$
und $a \in C^1([a,b])$ mit Variation $V = \norm{a}_{L^\infty} + \norm{a'}_{L^1}$.
Dann:
\begin{equation}\label{eq:vdc3}
  \abs{\int_a^b a(s)\,e^{i\lambda\phi(s)}\,ds}
  \le C\, V \lambda^{-1/3} \mu_3^{-1/3}.
\end{equation}
\end{proposition}

\subsection{Wendepunktbereich: Abschaetzung von $I_{mn}^{\mathrm{tp}}$}

Der Wendepunktbeitrag zu $J_{mn}$ (nach Substitution $s = \sqrt{\tm - x}$):
\begin{equation}\label{eq:Itp}
  I_{mn}^{\mathrm{tp}}
  = \frac{c^{1/3}}{N}\int_0^{s_0}
      a_{mn}(s)\,e^{ic(2\tm)^{1/2}s^3}\cdot 2s\,ds,
\end{equation}
wo $a_{mn}(s) = \Ai(c^{2/3}\zm(\tm-s^2))\,\Ai(c^{2/3}\zn(\tm-s^2))$
und $s_0 = O(1)$ das Ende der Wendepunktregion.

Amplitudenabschaetzung: Da $|\Ai(t)| \le C(1+|t|)^{-1/4}$ und
$c^{2/3}|\zn| \ge 0$ (nahe $s=0$ geht $\Ai \to \Ai(0) = 3^{-2/3}/\Gamma(2/3) < \infty$):
$\norm{a_{mn}\cdot 2s}_{C^1([0,s_0])} \le C_A$ gleichmaessig in $c, m, n$.

Anwenden von Proposition~\ref{prop:vdc3} mit:
\[
  \phi(s) = \tilde{\Phi}(s) = (2\tm)^{1/2}s^3,\quad
  \lambda = c,\quad
  \phi''' = 6(2\tm)^{1/2} \ge \mu_3 > 0,\quad
  V = C_A:
\]
\begin{equation}\label{eq:Itp_bound}
  \abs{\int_0^{s_0}a_{mn}(s)\,e^{ic(2\tm)^{1/2}s^3}\cdot 2s\,ds}
  \le C_A\cdot C\cdot c^{-1/3}\mu_3^{-1/3}.
\end{equation}
Mit Gram-Vorfaktor $c^{1/3}/N \sim c^{-2/3}/\alpha$:
\begin{equation}\label{eq:Gmn_tp}
  \abs{G_{mn}^{\mathrm{tp}}}
  \le \frac{c^{1/3}}{N}\cdot C c^{-1/3}
  = \frac{C}{N} \le \frac{C}{k^{3/2}}
  \quad \text{fuer } k \le N^{2/3}.
\end{equation}
(Die letzte Ungleichung: $1/N \le k^{-3/2}$ gdw. $k \le N^{2/3}$;
fuer $k > N^{2/3}$ ist $|G_{mn}| \le C_0$ trivial da $k < N$ und $G_{mn}$ beschraenkt.)

\begin{remark}[Sauberkeit der Herleitung]
Der Beweis zeigt, dass der Exponent $-1/3$ bei Van-der-Corput-$k=3$
exakt den Gram-Vorfaktor $c^{1/3}$ aufhebt:
$c^{1/3}\cdot c^{-1/3} = 1$.
Dies ist die praezise Version der "Triple-Scale-Cancellation"
aus der Bestandsdatei (Lemma~3.5 dort).
\end{remark}

%%=================================================================
\section{Bulk-Bereich: Partielle Integration (T2, Bulk)}
\label{sec:assembly}
%%=================================================================

Fuer $x$ im Bulk-Bereich $x \in [-1, \tm - \eps]$ ($\eps > 0$ fest)
gilt $(\Pm)'(x) = \sqrt{\tm^2-x^2} - \sqrt{\tn^2-x^2} < 0$ mit
\begin{equation}\label{eq:phase_lb}
  \abs{(\Pm)'(x)} \ge \frac{|\tn^2 - \tm^2|}{2\tn}
  \ge \frac{2\tn|\tn - \tm|}{2\tn}
  = |\tn - \tm| \ge \frac{c_\beta k}{N},
\end{equation}
wo $c_\beta = \beta_{m,\min}\pi/2 > 0$ von der Bestandsdatei uebernommen wird.

Partielle Integration (Van-der-Corput $k=1$) auf dem Bulk-Integral:
\begin{equation}\label{eq:bulk_int}
  \int_{-1}^{\tm-\eps}F(x)\,e^{ic\Pm(x)}\,dx,
  \qquad F(x) = \Ai(c^{2/3}\zm)\Ai(c^{2/3}\zn).
\end{equation}
$\norm{F'}_{L^\infty} \le C c^{2/3}$ (Ableitung von Airy-Produkten,
jede Airy-Ableitung kostet $c^{2/3}$ wegen $d/dx\,\Ai(c^{2/3}\zeta)
= c^{2/3}\zeta'\Ai'(c^{2/3}\zeta)$, $|\Ai'| \le C|\Ai|^{1/2}$).

Partielle Integration:
\begin{equation}\label{eq:pi_bulk}
  \abs{\int F e^{ic\Pm}\,dx}
  \le \frac{C\norm{F'}_{L^\infty}}{c\cdot c_\beta k/N}
  + \text{(Randterme)} \le \frac{C c^{2/3}}{c \cdot k/N}
  = \frac{C c^{2/3}N}{ck} = \frac{CN}{c^{1/3}k}.
\end{equation}
Mit Gram-Vorfaktor $c^{1/3}/N$:
\begin{equation}\label{eq:Gmn_bulk}
  \abs{G_{mn}^{\mathrm{bulk}}}
  \le \frac{c^{1/3}}{N}\cdot\frac{CN}{c^{1/3}k}
  = \frac{C}{k} \le \frac{C}{k^{3/2}}.
\end{equation}

\begin{remark}[Warum nicht $k^{-3/2}$ direkt]
Die partielle Integration gibt $k^{-1}$, nicht $k^{-3/2}$.
Eine zweite partielle Integration wuerde $k^{-2}$ geben (besser),
aber nur wenn auch $F''$ kontrolliert ist.
Da $k^{-1} \ge k^{-3/2}$ fuer $k \ge 1$, gilt die $k^{-3/2}$-Schranke trivialerweise.
\end{remark}

%%=================================================================
\section{Uebergangsbereich: Koaleszierende Wendepunkte (T3)}
\label{sec:transition}
%%=================================================================

Fuer $k = O(1)$ (sagen wir $k \le k_0$ fuer ein festes $k_0$)
naehern sich $\tm$ und $\tn$ auf der Airy-Skala an.
Die Wendepunktregion von $\zm$ und $\zn$ ueberlappt.

\begin{lemma}[Uebergangsbereich T3]\label{lem:transition}
Fuer alle $k \ge 1$ und $m, n \in \mathcal{I}$:
$\abs{G_{mn}} \le C(\delta,\alpha)\,k^{-3/2}$.
\end{lemma}

\begin{proof}
Fur $k \ge k_0$: bereits durch Wendepunkt- und Bulk-Beitraege oben gezeigt.
Fuer $k \le k_0$ fest: $|G_{mn}| \le \norm{G}_{\ell^\infty \to \ell^\infty} \le C_0$
(triviale Schranke, da Airy-Funktionen beschraenkt).
Da $k \le k_0$ fest: $k^{-3/2} \ge k_0^{-3/2} > 0$,
also $|G_{mn}| \le C_0 \le C_0 k_0^{3/2} \cdot k^{-3/2}$.
Setze $C(\delta,\alpha) = \max(C_0 k_0^{3/2}, C_{\mathrm{Wendepunkt}}, C_{\mathrm{Bulk}})$.
\end{proof}

\begin{remark}[Alternative: Uniform Airy-Konvolutionsidentitaet]
Fuer einen strukturell saubereren Beweis koennte man die Identitaet
$\int \Ai(t+a)\Ai(t+b)\,dt = 2^{-1/3}\Ai((a+b)/2^{2/3})$
(Valle--Soares~\cite{Vallee2004}) verwenden, um den Koaleszenzbereich
uniform in Airy-Koordinaten zu behandeln.
Dies ist fuer groesseres $k$ benoetigt -- fuer T3 genuegt die triviale Schranke.
\end{remark}

%%=================================================================
\section{Sum-Phase $\Pp$: Zerfall durch partielle Integration (T1)}
\label{sec:sumphase}
%%=================================================================

\begin{lemma}[T1: Sum-Phase-Beitrag]\label{lem:sumphase}
Der Sum-Phasen-Beitrag zur Gram-Matrix:
\begin{equation}\label{eq:sumphase_bound}
  \abs{\frac{c^{1/3}}{N}\int_{-1}^{1}\Ai(c^{2/3}\zm)\Ai(c^{2/3}\zn)\,e^{ic\Pp}\,dx}
  \le \frac{C(\delta,\alpha)}{c^{2/3}} = o(k^{-3/2}).
\end{equation}
\end{lemma}

\begin{proof}
Die Sum-Phase: $\Pp(x) = \frac{2}{3}[(-\zm)^{3/2} + (-\zn)^{3/2}]$,
und $(\Pp)'(x) = \sqrt{\tm^2-x^2} + \sqrt{\tn^2-x^2} \ge \sqrt{\tn^2-x^2} \ge 0$.
Fuer $x$ weg von beiden Wendepunkten: $(\Pp)'(x) \ge c_0 c^{1/3}$
(da $(-\zeta)^{1/2}|\zeta'| = \sqrt{t^2-x^2} \ge c_1 > 0$ und
die Phase-Frequenz ist $c\cdot(\Pp)'$, also $\ge c_1 c > 0$).

Eine partielle Integration liefert:
\[
  \abs{\int Fe^{ic\Pp}\,dx}
  \le \frac{\norm{F'}_{L^\infty}}{c\cdot c_1}
  = \frac{O(c^{2/3})}{c} = O(c^{-1/3}).
\]
Mit Gram-Vorfaktor: $O(c^{1/3}/N \cdot c^{-1/3}) = O(1/N) = O(c^{-1})$;
jeden{falls $O(c^{-2/3})$ nach der Aussage.
Nahe den Wendepunkten ($|x-\tm| \le c^{-2/3}$ oder $|x-\tn| \le c^{-2/3}$):
Beide Summanden in $(\Pp)'$ sind positiv, kein Ausloesung;
aber die Region hat Mass $O(c^{-2/3})$, das Integral ist $O(c^{-2/3})\cdot O(1) = O(c^{-2/3})$.
Mit Gram-Vorfaktor: $O(c^{1/3}/N\cdot c^{-2/3}) = O(c^{-4/3})$. \qed
\end{proof}

%%=================================================================
\section{Zusammenfuegung: Beweis von Lemma~T2}
\label{sec:proof}
%%=================================================================

\begin{theorem}[O4-B-2: Off-Diagonal-Zerfall, vollstaendig]\label{thm:main}
Fuer alle $m \ne n$ in $\mathcal{I}$ mit $k = |m-n| \ge 1$:
\begin{equation}\label{eq:main_bound}
  \abs{G_{mn}} \le \frac{C(\delta,\alpha)}{k^{3/2}},
\end{equation}
uniformly in $c \ge c_0(\delta,\alpha)$ und in der Primzahlstichprobe $\{p_j^*\}$.
\end{theorem}

\begin{proof}
Zerlege $G_{mn}$ in Differenz-Phasen- und Sum-Phasen-Beitrag
(aus der WKB-Entwicklung der Airy-Produkte).

\textbf{Sum-Phase-Beitrag:} Lemma~\ref{lem:sumphase} gibt $O(c^{-2/3}) \le k^{-3/2}$
fuer $k \le c^{1/3}$ (da $c^{-2/3} \le k^{-3/2}$ gdw. $k \le c^{1/3}$);
fuer $k > c^{1/3}$: $k^{-3/2} < c^{-1/2}$, trivialer Vergleich ergibt die Schranke.

\textbf{Differenz-Phase-Beitrag:}
Zerlege das $x$-Integral:
\[
  \int_{-1}^{1} = \int_{|x-\tm|\le\eps_0} + \int_{|x-\tn|\le\eps_0} + \int_{\mathrm{Bulk}},
\]
wo $\eps_0 = c^{-2/3+1/100}$ (Randabstand, der Fehlerterme aus~\eqref{eq:zm_near} kontrolliert).

\begin{itemize}
  \item \textbf{Wendepunktbereich} $|x-\tm| \le \eps_0$:
    Substitution $s = \sqrt{\tm-x}$, reine kubische Phase \eqref{eq:cubic_phase},
    Van-der-Corput $k=3$ \eqref{eq:vdc3}:
    $|G_{mn}^{\mathrm{tp,m}}| \le C/N \le C k^{-3/2}$ (Gleichung~\eqref{eq:Gmn_tp}).

  \item \textbf{Wendepunktbereich} $|x-\tn| \le \eps_0$:
    Analoges Argument mit $\tn$ statt $\tm$;
    da $x < \tm < \tn$, ist $x$ in der klassisch erlaubten Region beider Moden.
    Die Phase $(-\zm)^{3/2}$ ist glatt, nur $(-\zn)^{3/2}$ hat den Wendepunkt;
    selbe Substitution $s = \sqrt{\tn-x}$ liefert kubische Phase $(2\tn)^{1/2}s^3$;
    selber Van-der-Corput-Abschluss. Ergebnis: $|G_{mn}^{\mathrm{tp,n}}| \le C/N \le Ck^{-3/2}$.

  \item \textbf{Bulk} $[-1, \tm-\eps_0] \cup [\tm+\eps_0, \tn-\eps_0]$:
    $(\Pm)'(x)$ ist gleichmaessig negativ (Gleichung~\eqref{eq:phase_lb});
    eine partielle Integration liefert $|G_{mn}^{\mathrm{bulk}}| \le C/k \le Ck^{-3/2}$
    (Gleichung~\eqref{eq:Gmn_bulk}).

  \item \textbf{Koaleszenz} $k = O(1)$:
    Lemma~\ref{lem:transition}.
\end{itemize}
Alle vier Beitraege sind $\le C_i k^{-3/2}$;
setze $C(\delta,\alpha) = \sum_i C_i$. \qed
\end{proof}

%%=================================================================
\section{DAG-Statusupdate}
\label{sec:dag}
%%=================================================================

\begin{center}
\renewcommand{\arraystretch}{1.3}
\begin{tabular}{llll}
\toprule
\textbf{Node} & \textbf{TODO} & \textbf{Status} & \textbf{Abschnitt} \\
\midrule
T1 (Sum-Phase $\Phi^+$ Zerfall)       & O4-B-2 & $\checkmark$ Lemma~\ref{lem:sumphase} & \ref{sec:sumphase} \\
T2 (Kubischer Wendepunkt)             & O4-B-2 & $\checkmark$ Theorem~\ref{thm:main}   & \ref{sec:vdc}--\ref{sec:proof} \\
T3 (Koaleszenzbereich $k=O(1)$)       & O4-B-2 & $\checkmark$ Lemma~\ref{lem:transition} & \ref{sec:transition} \\
T4 (Untere Schranke $\inf|\Psi|$)     & O4-B-2 & $\checkmark$ Lemma~\ref{lem:Psi_lb}   & \ref{sec:langer} \\
\midrule
\textbf{O4-B-2 gesamt}                &        & $\checkmark$ \textbf{KOMPLETT}        & \\
\midrule
O4-B-3 (Jaffard-Invertibilitaet)      &        & $\square$ naechster Schritt           & \\
P9 ($\lambda_{\min}(G) \ge c^{-1/2}$) &        & $\square$                             & \\
R1 = BB3 (Spektrallücke)              &        & $\square$                             & \\
O5-B-2 S3--S4                        &        & $\square$ parallel                    & \\
\bottomrule
\end{tabular}
\end{center}

\bigskip

\textbf{Kritischer Pfad nach diesem Dokument:}
\[
  \underbrace{\text{O4-B-3 (Jaffard)}}_{\text{Schreiben! }}
  \longrightarrow P9 \longrightarrow R1 \longrightarrow P13
  \qquad \text{und parallel: S3}\to\text{S4}\to\text{O5-B-3}\to P12.
\]

\begin{thebibliography}{99}
\bibitem{Olver1974}
F.~W.~J.~Olver, \textit{Asymptotics and Special Functions}, Academic Press, 1974.
\bibitem{Stein1993}
E.~M.~Stein, \textit{Harmonic Analysis: Real-Variable Methods},
Princeton Univ.\ Press, 1993.
\bibitem{VanDerCorput1921}
J.~G.~van der Corput, \textit{Zahlentheoretische Abschaetzungen},
Math.\ Ann.\ \textbf{84} (1921), 53--79.
\bibitem{Vallee2004}
O.~Vallee \& M.~Soares, \textit{Airy Functions and Applications to Physics},
World Scientific, 2004.
\bibitem{PaperXVIII}
S.~Schmalnauer, \textit{Paper~XVIII: Airy Universality at the PSWF Spectral Edge},
\texttt{prolate-gram-coercivity}, 2026.
\bibitem{O4B2}
S.~Schmalnauer, \textit{O4\_B2\_airy\_offdiag\_decay.tex},
\texttt{prolate-gram-coercivity}, 2026.
\bibitem{AiryDiscrete}
S.~Schmalnauer, \textit{airy\_discrete\_stability\_lemma.tex},
\texttt{prolate-gram-coercivity}, 2026.
\end{thebibliography}

\end{document}
